% ==================================================
% USEPACKAGES VA THU VIEN LIEN QUAN
% Bien dich khuyen nghi: XeLaTeX
% ==================================================

% =========================
% FONT VA NGON NGU
% =========================
\usepackage{fontspec}
\usepackage[vietnamese]{babel}

% Giup go dau ngoac kep, gach noi, ligature tot hon voi XeLaTeX
\defaultfontfeatures{Ligatures=TeX}

% =========================
% TOAN HOC CO BAN
% Luu y:
% - Nap amsmath/mathtools TRUOC unicode-math
% - Khong nap amssymb khi da dung unicode-math
% =========================
\usepackage{amsmath}
\usepackage{amsthm}
\usepackage{mathtools}

% =========================
% FONT TOAN UNICODE
% unicode-math nen duoc nap SAU amsmath/mathtools
% =========================
\usepackage{unicode-math}

\IfFontExistsTF{Libertinus Serif}
{
    \setmainfont{Libertinus Serif}
    \setsansfont{Libertinus Sans}
    \setmathfont{Libertinus Math}
}
{
    \setmainfont{TeX Gyre Termes}
    \setsansfont{TeX Gyre Heros}
    \setmathfont{TeX Gyre Termes Math}
}

% =========================
% MAU SAC - HINH ANH - HINH VE
% =========================
\usepackage{xcolor}
\usepackage{graphicx}
\usepackage{tikz}
\usepackage{tikz-3dplot}

\usetikzlibrary{
  calc,
  angles,
  quotes,
  intersections,
  patterns,
  patterns.meta,
  arrows.meta,
  bending,
  positioning,
  decorations.pathreplacing,
  decorations.markings,
  shapes.geometric,
  fadings,
  3d,
  through
}

% =========================
% HINH HOC - BANG BIEN THIEN
% =========================
\usepackage{tkz-tab}
\usepackage{tkz-euclide}

% =========================
% DO THI HAM SO - PGFPLOTS
% =========================
\usepackage{pgfplots}
\pgfplotsset{compat=1.18}

\usepgfplotslibrary{fillbetween}
\usepgfplotslibrary{statistics}

% =========================
% BANG BIEU - COT - DANH SACH
% =========================
\usepackage{array}
\usepackage{tabularx}
\usepackage{booktabs}
\usepackage{multirow}
\usepackage{multicol}
\usepackage{tasks}
\usepackage{enumitem}
\usepackage{environ}

% =========================
% TRANG - HEADER - FOOTER - SPACING
% Chi nap package, cau hinh cu the de trong preamble-layout.tex
% =========================
\usepackage{geometry}
\usepackage{fancyhdr}
\usepackage{lastpage}
\usepackage{setspace}




% =========================
% HYPERLINK PDF
% Dat gan cuoi preamble
% =========================
\usepackage[
  unicode,
  colorlinks=true,
  linkcolor=blue,
  urlcolor=blue,
  citecolor=blue
]{hyperref}




\usepackage[most]{tcolorbox}
\usepackage{etoolbox}


